\makeabstract
{
Benchmarking the Coffea Framework at the Coffea-casa Analysis Facility
}
{
Fernanda Sophia Morais Laroca (1), Andrew Wightman (2), Kenneth Bloom (2), Oksana Shadura (2), Benjamin Tovar (3)
}
{
\textsuperscript{1}Amherst College, \textsuperscript{2}University of Nebraska-Lincoln, \textsuperscript{3}University of Notre Dame
}
{
The volume of data generated by the  High-Luminosity Large Hadron Collider (LHC) will be 100 times that of the current LHC, making dedicated computing facilities necessary for efficient data analysis. Coffea-casa is one such analysis facility, featuring an easy-to-use JupyterHub interface and access to scalable computing resources. Benchmarking the scalability of computing resources is essential to understand how to best utilize these resources for future physics analyses.
}
{
We run the full TopEFT analysis on 24 datasets, comprising over 80 million events, a total of 42 times. In each job, we primarily vary the number of workers to understand how this affects the analysis runtime. We also explore different data access protocols and versions of the coffea framework to understand their runtimes.
}
{
We find no significant differences in performance when accessing data via XRootD with or without the XCache configuration, even though XCache is optimized for faster access. This is likely because the accessed data was stored at the University of Nebraska-Lincoln. We find that the analysis runs five times slower in the updated version of the coffea framework, as compared to its previous usable version. 
}
{
--
}

\newpage
